\documentclass{situationist-preprint}

\title{The Reality of Dispositions}
\msid{MET-11-a}
\msfamily{MET-11}
\msdiscipline{Metaphysics}
\mstype{Principal manuscript}
\mskeywords{dispositions; powers; conditional analysis; finks; masks; categorical bases; intrinsic properties; laws of nature}

\begin{document}
\maketitle

\begin{abstract}
A glass is fragile, sugar is soluble, a magnet attracts iron. Are such dispositions real properties of their bearers, and if so, are they reducible to something else? Anti-realists hold that dispositional ascriptions are licences for inference rather than descriptions of properties. Reductive realists hold that dispositions are real but reducible, either to conditionals about what the bearer would do or to categorical properties together with laws. We argue that dispositions are real and irreducible properties of their bearers. They are real because they make a counterfactual difference to what happens, in the sense in which causes do. They are not reducible to conditionals, because finks, masks and mimics separate the disposition from every conditional that does not itself mention the disposition. And they are not reducible to categorical bases together with laws, by a dilemma: if laws are Humean summaries, the reduction makes intrinsic dispositions extrinsic; if laws govern, the reduction relocates irreducible modality in the laws instead of eliminating it. We conclude by showing that the thesis is consistent with a best-system account of laws.
\end{abstract}

\section{The problem}\label{s:problem}

Many of the properties we ascribe to things are characterized by what their bearers would do in certain circumstances. A fragile glass is one that would break if struck; a soluble lump of sugar is one that would dissolve if placed in water. Such properties are dispositions. Each is associated with a \emph{stimulus}, the circumstance in which it is manifested, and a \emph{manifestation}, what the bearer does in that circumstance. The question is whether dispositions are real properties of their bearers, and, if they are, whether they are anything over and above the non-dispositional properties of their bearers and the conditionals that describe what the bearers would do.

Three positions have dominated. The first is anti-realism. \citet{Ryle1949} held that to ascribe a disposition is not to report a state of the bearer but to issue an \enquote{inference ticket}, licensing the inference from the stimulus to the manifestation. There is no state of being fragile; there is only the truth of the hypothetical that the glass would break if struck. The second is reductive realism of the conditional kind. \citet{Carnap1936} attempted to introduce dispositional terms by reduction sentences, and the simple conditional analysis, on which $x$ is disposed to $M$ when $S$ just in case $x$ would $M$ if $S$, has been the default form of the view. The third is reductive realism of the categorical kind. On this view, dispositions are real, but each is identical to, or realized by, a categorical property of the bearer, its \emph{causal basis}, which in virtue of the laws of nature produces the manifestation in the stimulus conditions. \citet{PPJ1982} argued that every disposition has a causal basis, that dispositions are distinct from their bases, and that the bases rather than the dispositions do the causal work; \citet{Armstrong1997} held that the bases are categorical and that dispositional predicates refer to them under a description.

Against all three, a tradition that runs from \citet{Mellor1974} and \citet{Shoemaker1980} through \citet{Mumford1998}, \citet{Molnar2003}, \citet{Heil2003} and \citet{Bird2007} to \citet{Vetter2015} holds that dispositions are real and irreducible: that being fragile is a way the glass is, and that this way cannot be analysed in terms of what is not dispositional. We defend this conclusion, and we defend it without the further claim, made by some in that tradition, that all properties are dispositional.

\section{Framework}\label{s:framework}

\begin{definition}[Disposition]\label{d:disp}
A \emph{disposition} is a property $D$ such that its ascription specifies a stimulus $S$ and a manifestation $M$, and such that a bearer of $D$ is, in the absence of interference, one that would $M$ if $S$.
\end{definition}

The clause \enquote{in the absence of interference} is part of the characterization of a disposition, not an analysis of it. Its content is the subject of Section~\ref{s:conditionals}.

\begin{definition}[Reality]\label{d:real}
A property is \emph{real} if its instantiation by a bearer is a fact about the bearer which makes a difference to what happens, in the sense that some occurrences depend counterfactually on it, holding suitable other conditions fixed.
\end{definition}

\begin{definition}[Reduction]\label{d:reduce}
A disposition $D$ is \emph{reducible} if there are non-dispositional facts, together with truths stated without dispositional vocabulary, in which a bearer's having $D$ is wholly grounded, uniformly across its bearers.
\end{definition}

Definition~\ref{d:reduce} covers both reductive programmes: the conditional programme, on which the grounds are conditionals stated without dispositional vocabulary, and the categorical programme, on which they are categorical properties together with laws. It does not count as a reduction the grounding of one disposition in others together with structure: the fragility of a glass is plausibly grounded in the dispositions of its molecular bonds and their arrangement, and irreducibility in our sense allows this. What it excludes is that dispositionality should disappear at the end of the analysis.

\begin{principle}[Intrinsicness]\label{p:intr}
Typical dispositions are intrinsic: if two things are perfect duplicates, they have the same such dispositions.
\end{principle}

Principle~\ref{p:intr} was accepted by \citet{Lewis1997}, who used it to constrain the analysis of dispositions. \citet{McKitrick2003} has argued that some dispositions, such as being vulnerable or being recognizable, are extrinsic. We grant this; Principle~\ref{p:intr} concerns the dispositions, such as fragility, solubility and charge, that are the central cases.

\section{The argument}\label{s:argument}

\subsection{Reality}

\begin{proposition}[Reality]\label{p:reality}
Dispositions are real.
\end{proposition}

\begin{proof}
Consider two glasses struck in the same way, one fragile and one not, say because the second has been tempered. The first breaks and the second does not. Holding fixed the striking and the other conditions, whether the glass breaks depends on whether it is fragile. So the fragility of the first glass is a fact about it on which its breaking counterfactually depends, and, by Definition~\ref{d:real}, it is real. The argument generalizes to any disposition that is ever triggered in the absence of interference.
\end{proof}

On the account of causation of MET-04, a cause is a counterfactual difference-maker selected by norms of normality. The fragility of a glass is a difference-maker for its breaking; whether it is selected as a cause depends on the norms in play. In an ordinary context, where glasses are normally fragile, the striking is the cause and the fragility a standing condition, which is the familiar verdict. In a context where glasses are normally tempered, the fragility of this glass is a deviation, and is a cause. Proposition~\ref{p:reality} needs only the first component of the account: that dispositions make a difference, whether or not they are selected as causes.

Ryle's view cannot accept Proposition~\ref{p:reality}, because on it there is no state of the glass for the breaking to depend on; there is only the truth of a conditional. But conditionals require truthmakers, and what makes it true that this glass would break if struck, while the tempered glass would not, is a difference between the glasses. Ryle's inference tickets are valid of some objects and not of others, and the difference in their validity is a difference in the objects. To say that difference is the disposition is the realist's thesis.

\subsection{Against the conditional reduction}\label{s:conditionals}

\begin{lemma}[Finks, masks and mimics]\label{l:cond}
No conditional stated without dispositional vocabulary is necessary and sufficient for the possession of a disposition.
\end{lemma}

\begin{proof}
Consider the simple analysis: $x$ has $D$ if and only if $x$ would $M$ if $S$. It fails in both directions.

\emph{Finks} \citep{Martin1994}. A live wire is connected to an electro-fink, a device that detects when the wire is about to be touched and instantly makes it dead. The wire is live, but it is false that it would conduct if touched. Conversely, a dead wire with a reverse fink, which makes it live when touched, would conduct if touched, though it is dead. Finks separate the disposition from the conditional because the stimulus itself removes or confers the disposition.

\emph{Masks and antidotes} \citep{Johnston1992,Bird1998}. A fragile glass packed in protective material would not break if struck, though it remains fragile; a person who has taken a poison together with its antidote would not die, though the poison remains lethal. Masks separate the disposition from the conditional without removing the disposition: they interfere with the process from stimulus to manifestation.

\emph{Mimics}. A sturdy glass struck by a sorcerer who breaks whatever is struck would break if struck, but is not fragile. Mimics produce the manifestation without the disposition.

\citet{Lewis1997} proposed a reformed analysis designed for finks. On it, $x$ is disposed to $M$ when $S$ if and only if, for some intrinsic property $B$ that $x$ has, if $x$ were $S$ and retained $B$ for a sufficient time, then $S$ and $x$'s having $B$ would jointly be an $x$-complete cause of $x$'s $M$-ing. The analysis handles finks, since the fink removes $B$ and the counterfactual holds $B$ fixed. It does not handle masks, since the masked glass retains its intrinsic basis and the stimulus and basis are not, in the presence of the packing, a complete cause of breaking; nor does it handle antidotes, as \citet{Bird1998} argued. The only way to repair the analysis is to add a clause excluding interference, and interference is anything that would prevent the bearer's disposition from being manifested, which mentions the disposition. \citet{Choi2006} argued that the simple analysis survives if its stimulus is specified with care, so as to exclude masks and finks. But the specification that works is one that excludes every circumstance in which the disposition is prevented from being manifested, and a specification of that kind is not stated without dispositional vocabulary.

Every conditional that survives the three classes of counterexample therefore contains, explicitly or in its ceteris paribus clause, the notion of the disposition's being interfered with, and so is not stated without dispositional vocabulary.
\end{proof}

\citet{ManleyWasserman2008} drew a related conclusion from different cases and proposed that dispositions be linked to conditionals by proportion rather than by entailment: a thing is fragile if it would break in a suitable proportion of the cases of striking. The proposal is not offered as a reduction, and it is consistent with the present thesis. The companion manuscript MET-11-b examines finkish and masked dispositions in more detail.

\subsection{Against the categorical reduction}

\begin{lemma}[The dilemma]\label{l:cat}
No disposition that satisfies Principle~\ref{p:intr} is reducible to categorical properties together with laws.
\end{lemma}

\begin{proof}
Suppose that a disposition $D$ is wholly grounded in a categorical property $B$ of the bearer together with a law $L$ to the effect that $B$-things that are $S$ are $M$. Laws are either summaries of the course of events or something that governs it.

Suppose first that laws are summaries, as on the best system account defended in MET-12. Then whether $L$ is a law depends on the whole mosaic, including events arbitrarily far from the bearer: one $B$-thing elsewhere that is $S$ and not $M$ can make $L$ false, and an addition to the mosaic can change which generalizations belong to its best summary. Then a perfect duplicate of the bearer in a world with a different mosaic need not have $D$, although it has $B$, and $D$ is extrinsic, contrary to Principle~\ref{p:intr}.

Suppose instead that laws govern: that $L$ is, for instance, a relation of necessitation between the universals $B$ and $M$ \citep{Armstrong1983}. Then the ground of $D$ includes a relation whose nature is to make it the case that $B$-things that are $S$ are $M$. That relation is itself dispositional in the sense that matters, since its characterization specifies a stimulus and a manifestation, and the reduction has relocated the dispositional element rather than eliminated it \citep{Mumford2004}. The ground is therefore not wholly non-dispositional, and $D$ is not reduced in the sense of Definition~\ref{d:reduce}.

In either case, $D$ is not reducible to categorical properties together with laws.
\end{proof}

The second horn has a familiar history. \citet{vanFraassen1989} asked how a relation between universals could entail a regularity, and \citet{Lewis1983} observed that calling the relation necessitation does not give it the power to necessitate. The anti-Humean who answers that the relation necessitates by its nature has accepted a primitive whose nature is dispositional, which is the irreducibility we claim. The first horn is the one that bears on MET-12, and we return to it in Section~\ref{s:laws}.

\begin{theorem}\label{t:main}
Dispositions are real and irreducible properties of their bearers.
\end{theorem}

\begin{proof}
Reality is Proposition~\ref{p:reality}. A reduction of a disposition would ground it wholly in non-dispositional facts together with truths stated without dispositional vocabulary. Such truths are either conditionals about the bearer or laws, together with the categorical properties of the bearer. The first option is excluded by Lemma~\ref{l:cond}, the second by Lemma~\ref{l:cat}.
\end{proof}

The proof assumes that the two programmes exhaust the options for reduction. We know of no third. A reduction to the manifestations themselves, on which being fragile is breaking at some time, is excluded by the fact that unmanifested dispositions are possible: a fragile glass that is never struck, and is melted down, was fragile.

\section{Dispositions and laws}\label{s:laws}

Lemma~\ref{l:cat} uses the best system account of laws in one horn of a dilemma, and MET-12 defends that account. It may seem that the two manuscripts pull against each other: Lewis's Humeanism holds that the fundamental properties are categorical and that there are no necessary connections between distinct existences, while irreducible dispositions are precisely such connections.

The tension is with Lewis's package, not with the best system account of laws. As MET-12 argues, the account of what makes a generalization a law is independent of the thesis that the mosaic is categorical. If dispositions are irreducible, they are part of the mosaic, and the best system summarizes, among other things, the regularities of their manifestation. The laws so understood do not ground dispositions, since lawhood is global and dispositions are intrinsic; this is the first horn of Lemma~\ref{l:cat} read in the other direction. Nor do dispositions make the laws idle: the laws remain the best summary of the mosaic and explain its regularities by subsumption, while the dispositions do the producing, where anything does. The position lies between Lewis, for whom the mosaic is categorical and the laws summarize it, and \citet{Bird2007}, for whom the mosaic is dispositional and the laws are derived from the essences of its properties. It takes the irreducible dispositions from Bird and the account of lawhood from Lewis.

\section{Objections and replies}\label{s:objections}

\begin{objection}[Dormitive virtue]
To explain why opium puts people to sleep by its dormitive virtue is to explain nothing. Explanations by dispositions are empty, and a property that does no explanatory work is not real.
\end{objection}

\begin{reply}
The explanation of sleep by dormitive virtue is uninformative to someone who already knows that opium puts people to sleep, since it adds only that the effect is due to the drug. But that addition is not nothing: it locates the source of the effect in the drug rather than in the patient or the circumstances, and it contrasts opium with substances that would not have had the effect. Proposition~\ref{p:reality} claims no more than this. That a dispositional explanation is shallow does not make it false; it is completed, not replaced, by an account of the mechanism, which typically cites further dispositions of the drug's constituents.
\end{reply}

\begin{objection}[Causal impotence]
\citet{PPJ1982} argued that dispositions are causally impotent: whenever a fragile glass breaks, its causal basis, the molecular structure of the glass, is sufficient for the breaking, and the disposition adds nothing.
\end{objection}

\begin{reply}
The argument assumes that a sufficient condition excludes other causes at a different level of description. On a difference-making account it does not. Whether the glass breaks depends on whether it is fragile, and does not depend on the exact molecular structure, since many structures would have made the glass fragile. The disposition is the proportional difference-maker, and the basis is too specific. MND-11 makes the parallel point about mental causes, and the reasoning is the same. Moreover, on Lemma~\ref{l:cat} the basis does not by itself, without dispositions, ensure the breaking at all.
\end{reply}

\begin{objection}[Pointing to the non-existent]
A disposition is directed at its manifestation, which may never occur. A property essentially related to something that does not exist is a Meinongian absurdity \citep{Armstrong1997}.
\end{objection}

\begin{reply}
The disposition is related not to a particular manifestation event but to a type of manifestation, breaking, which is instantiated elsewhere if not here. \citet{Molnar2003} described the directedness of powers as a physical analogue of intentionality, and the analogy is apt in the respect that intentional states need not have existing particular objects. Nothing in Theorem~\ref{t:main} requires a relation to a non-existent particular.
\end{reply}

\begin{objection}[Regress]
If every property is a disposition to acquire further dispositions, nothing ever happens except the passing of powers, and the world is all potentiality and no act \citep{Armstrong1997}.
\end{objection}

\begin{reply}
The objection, whatever its force, targets pandispositionalism, the view that every property is dispositional. Theorem~\ref{t:main} does not assert that view. It asserts that dispositions are real and irreducible; it is consistent with there being categorical properties, including properties whose instances are the manifestations of dispositions.
\end{reply}

\section{What the result does not establish}\label{s:limits}

The result does not establish that every property is dispositional, or that the fundamental properties are. It does not settle whether dispositions are identical to qualities, as \citet{Heil2003} proposes, or whether the modality of dispositions is the source of all modality, as \citet{Vetter2015} argues. It does not provide a positive analysis of the ceteris paribus clause in Definition~\ref{d:disp}, and Lemma~\ref{l:cond} argues that none in non-dispositional terms is available. Finally, Lemma~\ref{l:cat} applies only to intrinsic dispositions; an extrinsic disposition may be partly grounded in the environment of its bearer, though not, if MET-12 is right, in the laws.

\section{Conclusion}

Dispositions are real, because what happens depends on them, and irreducible, because neither conditionals nor categorical bases with laws can do their work: conditionals are separated from dispositions by finks, masks and mimics, and the categorical reduction makes dispositions extrinsic if laws are Humean and relocates them in the laws if laws govern. The thesis is consistent with the best system account of laws, provided that account is not tied to the claim that the mosaic is categorical.

\begin{thebibliography}{99}\small
\bibitem[Armstrong(1983)]{Armstrong1983} Armstrong, D.~M. (1983). \emph{What Is a Law of Nature?} Cambridge University Press.
\bibitem[Armstrong(1997)]{Armstrong1997} Armstrong, D.~M. (1997). \emph{A World of States of Affairs}. Cambridge University Press.
\bibitem[Bird(1998)]{Bird1998} Bird, A. (1998). Dispositions and antidotes. \emph{Philosophical Quarterly}, 48(191), 227--234.
\bibitem[Bird(2007)]{Bird2007} Bird, A. (2007). \emph{Nature's Metaphysics: Laws and Properties}. Oxford University Press.
\bibitem[Carnap(1936)]{Carnap1936} Carnap, R. (1936). Testability and meaning. \emph{Philosophy of Science}, 3(4), 419--471.
\bibitem[Choi(2006)]{Choi2006} Choi, S. (2006). The simple vs. reformed conditional analysis of dispositions. \emph{Synthese}, 148(2), 369--379.
\bibitem[Heil(2003)]{Heil2003} Heil, J. (2003). \emph{From an Ontological Point of View}. Oxford University Press.
\bibitem[Johnston(1992)]{Johnston1992} Johnston, M. (1992). How to speak of the colors. \emph{Philosophical Studies}, 68(3), 221--263.
\bibitem[Lewis(1983)]{Lewis1983} Lewis, D. (1983). New work for a theory of universals. \emph{Australasian Journal of Philosophy}, 61(4), 343--377.
\bibitem[Lewis(1997)]{Lewis1997} Lewis, D. (1997). Finkish dispositions. \emph{Philosophical Quarterly}, 47(187), 143--158.
\bibitem[Manley and Wasserman(2008)]{ManleyWasserman2008} Manley, D. and Wasserman, R. (2008). On linking dispositions and conditionals. \emph{Mind}, 117(465), 59--84.
\bibitem[Martin(1994)]{Martin1994} Martin, C.~B. (1994). Dispositions and conditionals. \emph{Philosophical Quarterly}, 44(174), 1--8.
\bibitem[McKitrick(2003)]{McKitrick2003} McKitrick, J. (2003). A case for extrinsic dispositions. \emph{Australasian Journal of Philosophy}, 81(2), 155--174.
\bibitem[Mellor(1974)]{Mellor1974} Mellor, D.~H. (1974). In defense of dispositions. \emph{Philosophical Review}, 83(2), 157--181.
\bibitem[Molnar(2003)]{Molnar2003} Molnar, G. (2003). \emph{Powers: A Study in Metaphysics}. Oxford University Press.
\bibitem[Mumford(1998)]{Mumford1998} Mumford, S. (1998). \emph{Dispositions}. Oxford University Press.
\bibitem[Mumford(2004)]{Mumford2004} Mumford, S. (2004). \emph{Laws in Nature}. Routledge.
\bibitem[Prior et~al.(1982)]{PPJ1982} Prior, E.~W., Pargetter, R. and Jackson, F. (1982). Three theses about dispositions. \emph{American Philosophical Quarterly}, 19(3), 251--257.
\bibitem[Ryle(1949)]{Ryle1949} Ryle, G. (1949). \emph{The Concept of Mind}. Hutchinson.
\bibitem[Shoemaker(1980)]{Shoemaker1980} Shoemaker, S. (1980). Causality and properties. In P.~van Inwagen (ed.), \emph{Time and Cause}. Reidel.
\bibitem[van Fraassen(1989)]{vanFraassen1989} van Fraassen, B.~C. (1989). \emph{Laws and Symmetry}. Oxford University Press.
\bibitem[Vetter(2015)]{Vetter2015} Vetter, B. (2015). \emph{Potentiality: From Dispositions to Modality}. Oxford University Press.
\end{thebibliography}

\end{document}
